% Generated from lessons/0012-1-partial-derivatives-total-differential.html; edit the HTML source, then rerun the converter.
\section{偏导数与全微分}
\lessonmark{偏导数与全微分}

这一节是从“多元连续”进入“多元微分”的第一关。考试最常考的不是单纯求偏导，而是分清：偏导存在、方向导数存在、可微、连续偏导、梯度、高阶混合偏导之间到底谁能推出谁。

\subsection*{本节主线}

\begin{conceptbox}
\textbf{偏导数}
只让一个变量动，其余变量冻结。它是一维导数的局部切片。
\end{conceptbox}

\begin{conceptbox}
\textbf{全微分}
让所有变量同时动，看函数增量是否能被一个线性函数主导。
\end{conceptbox}

\begin{conceptbox}
\textbf{梯度与 Jacobi 矩阵}
实值函数用梯度描述最速变化；向量值函数用 Jacobi 矩阵统一表示导数。
\end{conceptbox}

\begin{notebox}
一句话记忆：偏导是“单方向切片”；可微是“附近整体像一个线性函数”。多元里，偏导存在远远不够。
\end{notebox}

\subsection*{一、偏导数：冻结其他变量}

设 \(f\) 在开集 \(D\subset R^2\) 上有定义，\((x_0,y_0)\in D\)。若极限

\[
      \lim_{\Delta x\to0}
      \frac{f(x_0+\Delta x,y_0)-f(x_0,y_0)}{\Delta x}
    \]

存在，就称它为 \(f\) 在 \((x_0,y_0)\) 关于 \(x\) 的偏导数，记作

\[
      f_x(x_0,y_0),\qquad \frac{\partial f}{\partial x}(x_0,y_0).
    \]

同理可定义 \(f_y(x_0,y_0)\)。对 \(n\) 元函数，关于第 \(i\) 个变量的偏导数是

\[
      \frac{\partial f}{\partial x_i}(x^0)=
      \lim_{\Delta x_i\to0}
      \frac{f(x_1^0,\ldots,x_i^0+\Delta x_i,\ldots,x_n^0)-f(x^0)}
      {\Delta x_i}.
    \]

几何上，\(f_x(x_0,y_0)\) 是曲面 \(z=f(x,y)\) 被平面 \(y=y_0\) 切出的曲线在 \(x\) 方向的斜率。

\subsubsection*{偏导存在不保证连续}

经典反例是

\[
      f(x,y)=
      \begin{cases}
      \dfrac{xy}{x^2+y^2}, & (x,y)\ne(0,0),\\
      0, & (x,y)=(0,0).
      \end{cases}
    \]

在原点 \(f_x(0,0)=f_y(0,0)=0\)，但沿 \(y=x\) 趋于原点时 \(f(x,x)=1/2\)，所以函数在原点不连续。

\subsection*{二、方向导数：沿任意方向靠近}

令单位方向向量 \(v=(\cos\alpha,\sin\alpha)\)。教材采用单侧方向导数：

\[
      \frac{\partial f}{\partial v}(x_0,y_0)=
      \lim_{t\to0+}
      \frac{f(x_0+t\cos\alpha,y_0+t\sin\alpha)-f(x_0,y_0)}{t}.
    \]

方向导数是“只沿射线走”的变化率；偏导数是坐标轴方向上的特殊方向导数。

\begin{notebox}
方向导数存在也不等于可微。它只检查很多条直线，而可微要控制所有趋近方式。
\end{notebox}

\subsection*{三、全微分：整体线性近似}

设

\[
      \Delta z=f(x_0+\Delta x,y_0+\Delta y)-f(x_0,y_0),\qquad
      \rho=\sqrt{(\Delta x)^2+(\Delta y)^2}.
    \]

若存在与 \(\Delta x,\Delta y\) 无关的常数 \(A,B\)，使

\[
      \Delta z=A\Delta x+B\Delta y+o(\rho),
      \qquad \rho\to0,
    \]

则称 \(f\) 在 \((x_0,y_0)\) 可微。线性主部

\[
      dz=A\,dx+B\,dy
    \]

称为全微分。若 \(f\) 可微，则 \(f\) 连续，且偏导数存在，并且

\[
      A=f_x(x_0,y_0),\qquad B=f_y(x_0,y_0),\qquad
      df=f_x\,dx+f_y\,dy.
    \]

\subsubsection*{考试最重要的推出关系}

\begin{center}
\small
\begin{tabularx}{\linewidth}{|>{\raggedright\arraybackslash}p{0.18\linewidth}|>{\raggedright\arraybackslash}p{0.42\linewidth}|>{\raggedright\arraybackslash}X|}
\hline
\textbf{条件} & \textbf{能推出什么} & \textbf{不能偷换成什么} \\ \hline
\(f_x,f_y\) 在邻域存在且在该点连续 & \(f\) 在该点可微 & 这只是充分条件，不是必要条件。 \\ \hline
\(f\) 可微 & \(f\) 连续，偏导存在，方向导数由梯度给出 & 不要求偏导在邻域连续。 \\ \hline
偏导存在 & 只能说明坐标方向差商存在 & 不能推出连续，也不能推出可微。 \\ \hline
所有方向导数存在 & 只能说明所有直线方向差商存在 & 不能推出连续，也不能推出可微。 \\ \hline
\end{tabularx}
\end{center}

\subsection*{四、三个核心定理}

\subsubsection*{定理 12.1.1：可微推出方向导数公式}

若 \(f\) 在 \((x_0,y_0)\) 可微，则对任意单位方向 \(v=(\cos\alpha,\sin\alpha)\)，方向导数存在，且

\[
      \frac{\partial f}{\partial v}(x_0,y_0)
      =f_x(x_0,y_0)\cos\alpha+f_y(x_0,y_0)\sin\alpha.
    \]

证明思路：把 \(\Delta x=t\cos\alpha,\ \Delta y=t\sin\alpha\) 代入可微定义，再除以 \(t\)，让 \(t\to0+\)。

\subsubsection*{定理 12.1.2：连续偏导是可微的充分条件}

若 \(f_x,f_y\) 在 \((x_0,y_0)\) 的某邻域存在，并且在 \((x_0,y_0)\) 连续，则 \(f\) 在该点可微。

证明思路是把总增量拆成两段：

\[
      \begin{aligned}
      &f(x_0+\Delta x,y_0+\Delta y)-f(x_0,y_0)\\
      ={}&[f(x_0+\Delta x,y_0+\Delta y)-f(x_0,y_0+\Delta y)]\\
      &+[f(x_0,y_0+\Delta y)-f(x_0,y_0)].
      \end{aligned}
    \]

分别对一元函数用 Lagrange 中值定理，再利用偏导连续性，把系数替换成 \(f_x(x_0,y_0),f_y(x_0,y_0)\)，剩下的误差就是 \(o(\rho)\)。

\subsubsection*{定理 12.1.3：混合偏导相等的条件}

如果二阶混合偏导 \(f_{xy},f_{yx}\) 在 \((x_0,y_0)\) 的邻域存在，并且在该点连续，则

\[
      f_{xy}(x_0,y_0)=f_{yx}(x_0,y_0).
    \]

证明思路：构造矩形差分

\[
      f(x_0+h,y_0+k)-f(x_0+h,y_0)-f(x_0,y_0+k)+f(x_0,y_0),
    \]

先按 \(x\) 再按 \(y\) 使用一元中值定理，得到一个混合偏导；换顺序再做一次，得到另一个混合偏导。最后令 \(h,k\to0\)，连续性迫使二者相等。

\subsection*{五、梯度：方向导数的压缩写法}

若 \(f\) 可微，定义梯度

\[
      \nabla f(x_0,y_0)=
      \left(f_x(x_0,y_0),f_y(x_0,y_0)\right).
    \]

方向导数公式可写成

\[
      \frac{\partial f}{\partial v}(x_0,y_0)=\nabla f(x_0,y_0)\cdot v.
    \]

因此当 \(\nabla f\ne0\) 时，函数增长最快的单位方向是

\[
      v=\frac{\nabla f}{|\nabla f|},
    \]

最大方向导数是 \(|\nabla f|\)，最小方向导数是 \(-|\nabla f|\)，对应方向是 \(-\nabla f/|\nabla f|\)。若 \(\nabla f=0\)，所有方向的一阶线性变化率都是 \(0\)，不能说有唯一“最快方向”。

\subsection*{六、高阶偏导与高阶微分}

二阶偏导包括

\[
      f_{xx},\quad f_{xy},\quad f_{yx},\quad f_{yy}.
    \]

混合偏导不总相等；只有在适当连续条件下才可以换序。若二阶偏导连续，则二阶全微分为

\[
      d^2z=z_{xx}\,dx^2+2z_{xy}\,dx\,dy+z_{yy}\,dy^2.
    \]

若 \(k\) 阶偏导连续，可用算子形式记忆：

\[
      d^kz=\left(dx\,\frac{\partial}{\partial x}+dy\,\frac{\partial}{\partial y}\right)^kz.
    \]

对 \(n\) 元函数 \(u=f(x_1,\ldots,x_n)\)，对应为

\[
      d^ku=
      \left(dx_1\frac{\partial}{\partial x_1}+\cdots+
      dx_n\frac{\partial}{\partial x_n}\right)^ku.
    \]

\subsection*{七、向量值函数与 Jacobi 矩阵}

设 \(f:D\subset R^n\to R^m\)，写成分量形式

\[
      f(x)=
      \begin{pmatrix}
      f_1(x_1,\ldots,x_n)\\
      \vdots\\
      f_m(x_1,\ldots,x_n)
      \end{pmatrix}.
    \]

若各分量对各变量可偏导，则 Jacobi 矩阵是

\[
      J_f(x)=
      \begin{pmatrix}
      \dfrac{\partial f_1}{\partial x_1} & \cdots & \dfrac{\partial f_1}{\partial x_n}\\
      \vdots & \ddots & \vdots\\
      \dfrac{\partial f_m}{\partial x_1} & \cdots & \dfrac{\partial f_m}{\partial x_n}
      \end{pmatrix}.
    \]

向量值函数可微时，微分统一写成

\[
      dy=f'(x^0)\,dx=J_f(x^0)\,dx.
    \]

定理 12.1.4 的核心：向量值函数 \(f\) 在 \(x^0\) 可微，当且仅当每个坐标分量函数 \(f_i\) 在 \(x^0\) 可微。

\subsection*{八、即时判断练习}

\begin{quickcheck}
若 \(f_x(0,0),f_y(0,0)\) 都存在，能否推出 \(f\) 在 \((0,0)\) 连续？

\tcblower
\textbf{答案：}不一定
\end{quickcheck}

\begin{quickcheck}
若 \(f_x,f_y\) 在邻域存在且在点 \((x_0,y_0)\) 连续，能否推出 \(f\) 在该点可微？

\tcblower
\textbf{答案：}能推出
\end{quickcheck}

\begin{quickcheck}
若 \(f\) 可微且 \(\nabla f(p)\ne0\)，函数在 \(p\) 点沿哪个单位方向增长最快？

\tcblower
\textbf{答案：}梯度方向
\end{quickcheck}

\subsection*{九、课后题训练}

下面按教材第 135-137 页习题 1-22 整理。题面完整保留，提示只给入口；写作业时仍建议先独立算一遍，再展开提示核对方向。

\begin{exercisebox}
\textbf{习题 1\badge{偏导计算}}

求下列函数的偏导数：

\[
      \begin{aligned}
      &(1)\ z=x^5-6x^4y^2+y^6; \qquad
      (2)\ z=x^2\ln(x^2+y^2);\\
      &(3)\ z=xy+\frac{x}{y}; \qquad
      (4)\ z=\sin(xy)+\cos^2(xy);\\
      &(5)\ z=e^x(\cos y+x\sin y); \qquad
      (6)\ z=\tan\left(\frac{x^2}{y}\right);\\
      &(7)\ z=\sin\frac{x}{y}\cdot\cos\frac{y}{x}; \qquad
      (8)\ z=(1+xy)^x;\\
      &(9)\ z=\ln(x+\ln y); \qquad
      (10)\ z=\arctan\frac{x+y}{1-xy};\\
      &(11)\ u=e^{x(x^2+y^2+z^2)}; \qquad
      (12)\ u=x^{1/z};\\
      &(13)\ u=\frac1{\sqrt{x^2+y^2+z^2}}; \qquad
      (14)\ u=x^{y^z};\\
      &(15)\ u=\sum_{i=1}^n a_i x_i\quad (a_i\text{ 为常数});\\
      &(16)\ u=\sum_{i,j=1}^n a_{ij}x_i x_j\quad (a_{ij}=a_{ji}\text{ 且为常数}).
      \end{aligned}
      \]

\begin{hintbox}
含变量指数时先取对数：\(x^{1/z}=\exp((1/z)\ln x)\)，\(x^{y^z}=\exp(y^z\ln x)\)。二次型第 (16) 小问利用对称性，\(\partial u/\partial x_k=2\sum_j a_{kj}x_j\)。
\end{hintbox}
\end{exercisebox}

\begin{exercisebox}
\textbf{习题 2\badge{指定点偏导}}

设 \(f(x,y)=x+y-\sqrt{x^2+y^2}\)，求 \(f_x(3,4)\) 及 \(f_y(3,4)\)。

\begin{hintbox}
先求一般偏导，再代入 \((3,4)\)。这里 \(\sqrt{3^2+4^2}=5\)。
\end{hintbox}
\end{exercisebox}

\begin{exercisebox}
\textbf{习题 3\badge{偏微分方程验证}}

设 \(z=e^{x/y^2}\)，验证

\[
        2x\frac{\partial z}{\partial x}
        +y\frac{\partial z}{\partial y}=0.
      \]

\begin{hintbox}
把指数 \(x/y^2\) 看作复合函数，分别求 \(z_x,z_y\)，两项会抵消。
\end{hintbox}
\end{exercisebox}

\begin{exercisebox}
\textbf{习题 4\badge{空间曲线切线}}

曲线

\[
        \begin{cases}
        z=\dfrac{x^2+y^2}{4},\\
        y=4
        \end{cases}
      \]

在点 \((2,4,5)\) 处的切线与 \(x\) 轴的正向所夹的角度是多少？

\begin{hintbox}
令 \(y=4\)，把曲线参数化为 \((x,4,(x^2+16)/4)\)。在 \(x=2\) 处求切向量，再用点积求夹角。
\end{hintbox}
\end{exercisebox}

\begin{exercisebox}
\textbf{习题 5\badge{指定点全微分}}

求下列函数在指定点的全微分：

\[
      \begin{aligned}
      &(1)\ f(x,y)=3x^2y-xy^2,\quad (1,2);\\
      &(2)\ f(x,y)=\ln(1+x^2+y^2),\quad (2,4);\\
      &(3)\ f(x,y)=\frac{\sin x}{y^2},\quad (0,1)\text{ 和 }\left(\frac{\pi}{4},2\right).
      \end{aligned}
      \]

\begin{hintbox}
模板：\(df=f_x(a,b)\,dx+f_y(a,b)\,dy\)。
\end{hintbox}
\end{exercisebox}

\begin{exercisebox}
\textbf{习题 6\badge{全微分公式}}

求下列函数的全微分：

\[
      \begin{aligned}
      &(1)\ z=y^x; \qquad
      (2)\ z=xye^{xy};\\
      &(3)\ z=\frac{x+y}{x-y}; \qquad
      (4)\ z=\frac{y}{\sqrt{x^2+y^2}};\\
      &(5)\ u=\sqrt{x^2+y^2+z^2}; \qquad
      (6)\ u=\ln(x^2+y^2+z^2).
      \end{aligned}
      \]

\begin{hintbox}
对二元函数写 \(dz=z_xdx+z_ydy\)，对三元函数写 \(du=u_xdx+u_ydy+u_zdz\)。
\end{hintbox}
\end{exercisebox}

\begin{exercisebox}
\textbf{习题 7\badge{方向导数}}

求函数 \(z=xe^{2y}\) 在点 \(P(1,0)\) 处的沿从点 \(P(1,0)\) 到点 \(Q(2,-1)\) 方向的方向导数。

\begin{hintbox}
方向向量为 \(Q-P=(1,-1)\)，先单位化，再用 \(\nabla z(P)\cdot v\)。
\end{hintbox}
\end{exercisebox}

\begin{exercisebox}
\textbf{习题 8\badge{方向导数最值}}

设 \(z=x^2-xy+y^2\)，求它在点 \((1,1)\) 处沿方向 \(v=(\cos\alpha,\sin\alpha)\) 的方向导数，并指出：

\begin{enumerate}
 \item 沿哪个方向的方向导数最大？
 \item 沿哪个方向的方向导数最小？
 \item 沿哪个方向的方向导数为零？
 \end{enumerate}

\begin{hintbox}
先求 \(\nabla z(1,1)\)。最大方向是梯度方向，最小方向是反梯度方向，零方向与梯度垂直。
\end{hintbox}
\end{exercisebox}

\begin{exercisebox}
\textbf{习题 9\badge{由方向导数反推梯度}}

如果可微函数 \(f(x,y)\) 在点 \((1,2)\) 处的从点 \((1,2)\) 到点 \((2,2)\) 方向的方向导数为 \(2\)，从点 \((1,2)\) 到点 \((1,1)\) 方向的方向导数为 \(-2\)。求：

\begin{enumerate}
 \item 这个函数在点 \((1,2)\) 处的梯度；
 \item 点 \((1,2)\) 处的从点 \((1,2)\) 到点 \((4,6)\) 方向的方向导数。
 \end{enumerate}

\begin{hintbox}
前两个方向分别是 \((1,0)\) 和 \((0,-1)\)，所以可直接读出梯度的两个分量。
\end{hintbox}
\end{exercisebox}

\begin{exercisebox}
\textbf{习题 10\badge{求梯度}}

求下列函数的梯度：

\[
      \begin{aligned}
      &(1)\ z=x^2+y^2\sin(xy);\\
      &(2)\ z=1-\left(\frac{x^2}{a^2}+\frac{y^2}{b^2}\right);\\
      &(3)\ u=x^2+2y^2+3z^2+3xy+4yz+6x-2y-5z,\quad (1,1,1).
      \end{aligned}
      \]

\begin{hintbox}
梯度就是全部一阶偏导组成的向量。第 (3) 小问需要最后代入指定点。
\end{hintbox}
\end{exercisebox}

\begin{exercisebox}
\textbf{习题 11\badge{最速增加方向}}

对于函数 \(f(x,y)=xy\)，在第 I 象限（包括边界）的每一点，指出函数值增加最快的方向。

\begin{hintbox}
\(\nabla f=(y,x)\)。非原点处最快方向是 \((y,x)/\sqrt{x^2+y^2}\)。原点处梯度为零，没有唯一最快方向。
\end{hintbox}
\end{exercisebox}

\begin{exercisebox}
\textbf{习题 12\badge{偏导存在但方向导数异常}}

验证函数

\[
        f(x,y)=\sqrt[3]{xy}
      \]

在原点 \((0,0)\) 连续且可偏导，但除方向 \(e_i\) 和 \(-e_i\ (i=1,2)\) 外，在原点的沿其他方向的方向导数都不存在。

\begin{hintbox}
沿方向 \(v=(a,b)\) 计算 \(\frac{f(ta,tb)-f(0,0)}{t}\)。若 \(ab\ne0\)，会出现 \(t^{-1/3}\)。
\end{hintbox}
\end{exercisebox}

\begin{exercisebox}
\textbf{习题 13\badge{连续可偏导但不可微}}

验证函数

\[
      f(x,y)=
      \begin{cases}
      \dfrac{xy}{\sqrt{x^2+y^2}}, & x^2+y^2\ne0,\\
      0, & x^2+y^2=0
      \end{cases}
      \]

在原点 \((0,0)\) 连续且可偏导，但它在该点不可微。

\begin{hintbox}
原点偏导都为 \(0\)。若可微，应该有 \(f(x,y)=o(\sqrt{x^2+y^2})\)。沿 \(y=x\) 检查这个比值。
\end{hintbox}
\end{exercisebox}

\begin{exercisebox}
\textbf{习题 14\badge{偏导不连续但可微}}

验证函数

\[
      f(x,y)=
      \begin{cases}
      (x^2+y^2)\sin\dfrac1{x^2+y^2}, & x^2+y^2\ne0,\\
      0, & x^2+y^2=0
      \end{cases}
      \]

的偏导函数 \(f_x(x,y),f_y(x,y)\) 在原点 \((0,0)\) 不连续，但它在该点可微。

\begin{hintbox}
可微性用 \(|f(x,y)|\le x^2+y^2=o(\sqrt{x^2+y^2})\)。偏导不连续可从非原点偏导中出现的 \(\cos\frac1{x^2+y^2}\) 项看出。
\end{hintbox}
\end{exercisebox}

\begin{exercisebox}
\textbf{习题 15\badge{方向导数不够强}}

证明函数

\[
      f(x,y)=
      \begin{cases}
      \dfrac{2xy^2}{x^2+y^4}, & x^2+y^2\ne0,\\
      0, & x^2+y^2=0
      \end{cases}
      \]

在原点 \((0,0)\) 处沿各个方向的方向导数都存在，但它在该点不连续，因而不可微。

\begin{hintbox}
方向导数：令 \((x,y)=t(a,b)\)。不连续：沿抛物线 \(x=y^2\) 趋于原点。
\end{hintbox}
\end{exercisebox}

\begin{exercisebox}
\textbf{习题 16\badge{高阶导数}}

计算下列函数的高阶导数：

\[
      \begin{aligned}
      &(1)\ z=\arctan\frac{y}{x},\quad
      \frac{\partial^2z}{\partial x^2},\
      \frac{\partial^2z}{\partial x\partial y},\
      \frac{\partial^2z}{\partial y^2};\\
      &(2)\ z=x\sin(x+y)+y\cos(x+y),\quad
      \frac{\partial^2z}{\partial x^2},\
      \frac{\partial^2z}{\partial x\partial y},\
      \frac{\partial^2z}{\partial y^2};\\
      &(3)\ z=xe^{x/y},\quad
      \frac{\partial^3z}{\partial x^2\partial y},\
      \frac{\partial^3z}{\partial x\partial y^2};\\
      &(4)\ u=\ln(ax+by+cz),\quad
      \frac{\partial^4u}{\partial x^4},\
      \frac{\partial^4u}{\partial x^2\partial y^2};\\
      &(5)\ z=(x-a)^p(y-b)^q,\quad
      \frac{\partial^{p+q}z}{\partial x^p\partial y^q};\\
      &(6)\ u=xyze^{x+y+z},\quad
      \frac{\partial^{p+q+r}u}{\partial x^p\partial y^q\partial z^r}.
      \end{aligned}
      \]

\begin{hintbox}
连续高阶偏导可换序。第 (5) 小问会得到 \(p!q!\)。第 (6) 小问用 Leibniz 公式，比硬算更稳。
\end{hintbox}
\end{exercisebox}

\begin{exercisebox}
\textbf{习题 17\badge{高阶微分}}

计算下列函数的高阶微分：

\[
      \begin{aligned}
      &(1)\ z=x\ln(xy),\quad \text{求}\ d^2z;\\
      &(2)\ z=\sin^2(ax+by),\quad \text{求}\ d^3z;\\
      &(3)\ u=e^{x+y+z}(x^2+y^2+z^2),\quad \text{求}\ d^3u;\\
      &(4)\ z=e^x\sin y,\quad \text{求}\ d^kz.
      \end{aligned}
      \]

\begin{hintbox}
用算子法：\(d^k=(dx\,\partial_x+dy\,\partial_y+\cdots)^k\)。第 (4) 小问可利用 \(e^x\sin y=\operatorname{Im}(e^{x+iy})\)。
\end{hintbox}
\end{exercisebox}

\begin{exercisebox}
\textbf{习题 18\badge{由偏导反求函数}}

函数 \(z=f(x,y)\) 满足

\[
        \frac{\partial z}{\partial x}
        =-\sin y+\frac1{1-xy},
        \qquad
        f(0,y)=2\sin y+y^3.
      \]

求 \(f(x,y)\) 的表达式。

\begin{hintbox}
对 \(x\) 积分，积分“常数”应写成 \(C(y)\)。再用 \(f(0,y)\) 确定 \(C(y)\)。
\end{hintbox}
\end{exercisebox}

\begin{exercisebox}
\textbf{习题 19\badge{方程验证}}

验证：

\begin{enumerate}
 \item \(z=e^{-kn^2t}\sin(ny)\) 满足热传导方程 \[
          \frac{\partial z}{\partial t}=k\frac{\partial^2z}{\partial y^2}.
          \]
 \item \(u=z\arctan\frac{x}{y}\) 满足 Laplace 方程 \[
          \frac{\partial^2u}{\partial x^2}
          +\frac{\partial^2u}{\partial y^2}
          +\frac{\partial^2u}{\partial z^2}=0.
          \]
 \end{enumerate}

\begin{hintbox}
这类题不要解释物理意义，直接求导代入。第二小问注意 \(u\) 对 \(z\) 是一次的，所以 \(u_{zz}=0\)。
\end{hintbox}
\end{exercisebox}

\begin{exercisebox}
\textbf{习题 20\badge{参数确定}}

设

\[
        f(r,t)=t^\alpha e^{-r^2/(4t)},
      \]

确定 \(\alpha\)，使得 \(f\) 满足方程

\[
        \frac{\partial f}{\partial t}
        =
        \frac1{r^2}\frac{\partial}{\partial r}
        \left(r^2\frac{\partial f}{\partial r}\right).
      \]

\begin{hintbox}
两边都提出公共因子 \(f\)。左边为 \(f(\alpha/t+r^2/(4t^2))\)，右边会出现 \(f(-3/(2t)+r^2/(4t^2))\)，所以 \(\alpha=-3/2\)。
\end{hintbox}
\end{exercisebox}

\begin{exercisebox}
\textbf{习题 21\badge{向量值函数导数}}

求下列向量值函数在指定点的导数：

\[
      \begin{aligned}
      &(1)\ f(x)=(a\cos x,b\sin x,cx)^T,\quad x=\frac{\pi}{4};\\
      &(2)\ f(x,y,z)=(3x+e^y\cot z,\ x^3+y^3\tan z)^T,\quad
      (x,y,z)=\left(1,2,\frac{\pi}{4}\right);\\
      &(3)\ g(u,v)=(u\cos v,u\sin v,v)^T,\quad (u,v)=(1,\pi).
      \end{aligned}
      \]

\begin{hintbox}
把每个分量函数放在一行，每个自变量放在一列，组成 Jacobi 矩阵。
\end{hintbox}
\end{exercisebox}

\begin{exercisebox}
\textbf{习题 22\badge{Jacobi 矩阵反推分量}}

设 \(f:R^3\to R^3\) 为向量值函数。

\begin{enumerate}
 \item 如果坐标分量函数 \(f_1(x,y,z)=x,\ f_2(x,y,z)=y,\ f_3(x,y,z)=z\)，证明 \(f\) 的导数是单位阵；
 \item 写出坐标分量函数的一般形式，使 \(f\) 的导数是单位阵；
 \item 如果已知 \(f\) 的导数是对角阵 \(\operatorname{diag}(p(x),q(y),r(z))\)，那么坐标分量函数应该具有什么样的形式？
 \end{enumerate}

\begin{hintbox}
导数是单位阵表示 \(\partial f_i/\partial x_j=\delta_{ij}\)。导数是 \(\operatorname{diag}(p(x),q(y),r(z))\) 表示三个分量分别只由一个变量控制，再逐个积分。
\end{hintbox}
\end{exercisebox}

来源：陈纪修、于崇华、金路《数学分析》第三版下册，第十二章 §1，教材第 120-137 页。下一节：第十二章 §2：多元复合函数的求导法则。

上一节：第十一章 §3：连续函数的性质。复盘：错题集与好题集。
